\documentclass[11pt]{article}
\usepackage[margin=1in]{geometry}
\usepackage[utf8]{inputenc}
\usepackage[T1]{fontenc}
\begin{document}

\section*{Adaptive VR Rehabilitation System}

\subsection*{Revised Task Distribution and System Architecture Proposal}


\hrule


\subsection*{1. Introduction}


This document proposes an updated division of work and system architecture for the development of the \textbf{Adaptive VR Rehabilitation System for Upper-Limb Recovery Post-Stroke}.


The revision improves:

\begin{itemize}
\item workload balance between team members
\item alignment with individual academic tracks
\item technical depth of the project
\item inclusion of a data-driven adaptive component
\end{itemize}

\hrule


\subsection*{2. System Overview}


The system is designed as a \textbf{closed-loop adaptive rehabilitation platform}, where user performance continuously influences task difficulty.


\subsubsection*{Core Workflow}


User performs task -> movement is tracked -> trajectory is processed -> performance features are extracted -> data is stored -> adaptive model evaluates performance -> system adjusts difficulty


\hrule


\subsection*{3. Key Design Principles}


\begin{itemize}
\item Separation between \textbf{raw data acquisition} and \textbf{data interpretation}
\item Use of \textbf{standardized performance metrics across tasks}
\item Adoption of a \textbf{data-driven adaptive system}, extending the original rule-based approach
\item Modular architecture: Unity (interaction \& signal processing) + Backend (data \& intelligence)
\end{itemize}

\hrule


\subsection*{4. Revised Division of Work}


\subsubsection*{4.1 Shared Responsibilities}


Both students collaborate on:


\begin{itemize}
\item Design and implementation of VR environments:
\begin{itemize}
\item Kitchen scene
\item Living room scene
\item Garden scene
\end{itemize}
\end{itemize}

\begin{itemize}
\item Development of rehabilitation tasks:
\begin{itemize}
\item Watering plants
\item Picking objects
\item Hanging tools
\item Cleaning / manipulation tasks
\end{itemize}
\end{itemize}

\begin{itemize}
\item Definition of difficulty parameters:
\begin{itemize}
\item target size
\item reach distance
\item task complexity
\item visual guidance
\end{itemize}
\end{itemize}

\begin{itemize}
\item System testing and usability validation
\end{itemize}

\hrule


\subsubsection*{4.2 Mechatronics / Software (Student A)}


Responsible for \textbf{interaction systems and raw data acquisition}


\paragraph*{Motion Tracking}

\begin{itemize}
\item Capture hand/controller position and movement (3D coordinates)
\item Store trajectory temporarily during task execution
\item Manage tracking lifecycle:
\begin{itemize}
\item start tracking
\item update tracking
\item stop tracking
\end{itemize}
\end{itemize}

\paragraph*{Interaction System}

\begin{itemize}
\item Configure and standardize \textbf{SDK-based} reusable interaction components across tasks:
\begin{itemize}
\item grab / release
\item place (snapping, validation)
\item pour (thresholds, spill logic)
\item cut
\item task-specific rules, edge cases, and telemetry around SDK events
\end{itemize}
\end{itemize}

\paragraph*{Event Detection}

\begin{itemize}
\item Detect and log:
\begin{itemize}
\item task start / completion
\item errors
\item prompts / guidance usage
\end{itemize}
\end{itemize}

\paragraph*{Unity Integration}

\begin{itemize}
\item Integrate tracking and interaction systems within VR scenes
\item Ensure stable and responsive real-time behavior
\end{itemize}

\hrule


\subsubsection*{4.3 AI \& Data Analysis (Student B)}


Responsible for \textbf{performance modeling, data pipeline, and adaptive intelligence}


\paragraph*{Performance Modeling (Core Contribution)}

Define quantitative metrics derived from movement data:


\begin{itemize}
\item normalized completion time (task + difficulty based)
\item error count
\item prompt usage
\item path efficiency
\item movement smoothness
\item hesitation / pauses
\item accuracy
\end{itemize}

\paragraph*{Feature Extraction Logic}

\begin{itemize}
\item Design mathematical formulas for all metrics
\item Collaborate on implementation within Unity
\item Ensure consistency across tasks and sessions
\end{itemize}

\paragraph*{Data Architecture}

\begin{itemize}
\item Design relational database schema (PostgreSQL / Supabase)
\item Structure entities:
\begin{itemize}
\item patients (anonymized)
\item sessions
\item tasks
\item performance metrics
\item adaptation decisions
\end{itemize}
\end{itemize}

\paragraph*{Backend System}

\begin{itemize}
\item Develop API for:
\begin{itemize}
\item receiving performance data from Unity
\item storing structured data
\item returning adaptive decisions
\end{itemize}
\end{itemize}

\paragraph*{Adaptive System}


\textbf{Phase 1: Rule-Based Baseline}

\begin{itemize}
\item Define initial decision rules:
\begin{itemize}
\item high performance -> increase difficulty
\item moderate -> maintain
\item low -> decrease
\end{itemize}
\end{itemize}

\textbf{Phase 2: Data-Driven Model}

\begin{itemize}
\item Train machine learning model on collected data
\item Input: performance features
\item Output: difficulty adjustment decision
\end{itemize}

\paragraph*{Data Analysis \& Evaluation}

\begin{itemize}
\item Analyze user performance trends
\item Compare:
\begin{itemize}
\item rule-based vs machine learning approach
\end{itemize}
\item Evaluate system effectiveness and adaptability
\end{itemize}

\hrule


\subsection*{5. Movement Tracking and Feature Extraction Methodology}


The system performs \textbf{trajectory-based motion analysis} to derive meaningful performance metrics from raw movement data.


\subsubsection*{5.1 Trajectory Acquisition}


\begin{itemize}
\item Raw 3D positions are recorded during task execution
\item Data is stored temporarily in memory within Unity
\item No per-frame data is directly stored in the database
\end{itemize}

\hrule


\subsubsection*{5.2 Trajectory Preprocessing}


To improve signal quality:


\begin{itemize}
\item \textbf{Downsampling} is applied to reduce redundant points
\item \textbf{Smoothing techniques} (e.g., moving average or low-pass filtering) are used to reduce noise and jitter
\end{itemize}

\hrule


\subsubsection*{5.3 Ideal Path Definition}


For each task, an expected trajectory is defined:


\begin{itemize}
\item simple tasks -> straight-line path from start to target
\item multi-step tasks -> sequence of target checkpoints
\end{itemize}

\hrule


\subsubsection*{5.4 Feature Computation}


The following metrics are computed from the processed trajectory:


\paragraph*{Path Efficiency}

Ratio between optimal path and actual path:


efficiency = straight-line distance / actual path length


\hrule


\paragraph*{Deviation from Ideal Path}

\begin{itemize}
\item Average distance between actual trajectory and ideal path
\item Measures spatial accuracy of movement
\end{itemize}

\hrule


\paragraph*{Movement Smoothness}

Derived from kinematic properties:


\begin{itemize}
\item velocity variation
\item acceleration changes
\item optional jerk approximation
\end{itemize}

Smooth movement corresponds to low variability in motion


\hrule


\paragraph*{Speed Metrics}

\begin{itemize}
\item average speed
\item peak speed
\end{itemize}

\hrule


\paragraph*{Hesitation Detection}

\begin{itemize}
\item number and duration of pauses
\item based on velocity thresholds
\end{itemize}

\hrule


\subsubsection*{5.5 Data Storage Strategy}


Only \textbf{aggregated features} are stored in the database:


\begin{itemize}
\item completion time
\item efficiency
\item smoothness
\item errors
\item prompts
\item accuracy
\item hesitation (pause count, total pause duration)
\end{itemize}

Optional:

\begin{itemize}
\item raw trajectory may be stored separately (e.g., JSON or file storage) for offline analysis
\end{itemize}

\hrule


\subsubsection*{5.6 Rationale}


This approach:


\begin{itemize}
\item avoids storing large volumes of raw coordinate data
\item ensures efficient database usage
\item provides ML-ready structured features
\item aligns with signal processing and motion analysis principles
\end{itemize}

\hrule


\subsection*{6. System Architecture}


\subsubsection*{Unity (VR Layer)}

\begin{itemize}
\item task execution
\item motion tracking
\item trajectory preprocessing
\item feature computation
\end{itemize}

down


\subsubsection*{Backend (API Layer)}

\begin{itemize}
\item receives feature data
\item stores data in database
\item runs adaptive model inference
\end{itemize}

down


\subsubsection*{Database (PostgreSQL / Supabase)}

\begin{itemize}
\item stores structured performance data
\end{itemize}

down


\subsubsection*{Machine Learning Pipeline (Offline)}

\begin{itemize}
\item model training using collected dataset
\item deployment for runtime inference
\end{itemize}

\hrule


\subsection*{7. Adaptive Decision Process}


At runtime:


\begin{enumerate}
\item Unity computes performance features
\item Features are sent to backend
\item Backend evaluates:
\begin{itemize}
\item rule-based logic or trained model
\end{itemize}
\item Backend returns decision:
\begin{itemize}
\item increase difficulty
\item maintain
\item decrease
\end{itemize}
\item Unity updates task parameters accordingly
\end{enumerate}

\hrule


\subsection*{8. Improvements Over Original Proposal}


Compared to the initial project definition:


\begin{itemize}
\item VR development is now \textbf{shared}, improving workload balance
\item The tracking system is extended into a \textbf{full motion analysis pipeline}
\item The adaptive system evolves from \textbf{rule-based to data-driven}
\item A complete \textbf{machine learning component} is introduced
\item The system becomes a \textbf{closed-loop adaptive system}
\end{itemize}

\hrule


\subsection*{9. Work Requirements}


\begin{itemize}
\item Organized file structure in GitHub and Git LFS
\item Agreed versioning across plugins and applications
\item Efficient organization and consistency across previous and newly developed systems
\end{itemize}

\subsection*{10. Conclusion}


This revised structure ensures:


\begin{itemize}
\item balanced distribution of workload
\item clear separation of responsibilities
\item strong alignment with academic tracks
\item enhanced technical and research value
\end{itemize}

The project evolves from a VR application with tracking into a \textbf{data-driven adaptive rehabilitation system}, integrating motion analysis, performance modeling, and intelligent difficulty adaptation.




\end{document}


